两个输入必须具有相同的秘密密钥、环、模数 $q$ 和尺度 $\Delta=q/t$；所述模 $t$ 的明文运算假设 $t\mid q$。在这些运算章节中，$\textsf{GLWE}_{S,\sigma}(\text{phase})$ 是具有所列相位的密文的简写；求值后的密文不一定具有新鲜加密时相同的高斯分布或标准差 $\sigma$。所有噪声界均针对系数约减前的整数提升。

\textbf{- 参考资料：}
\href{https://www.zama.ai/post/tfhe-deep-dive-part-2}{深入理解 TFHE（二）：编码与线性分层运算}~\cite{tfhe-2}


$ $

假设有两个 GLWE 密文，分别加密不同的明文 $M^{\langle 1 \rangle}, M^{\langle 2 \rangle}$：

$\textsf{GLWE}_{S, \sigma}(\Delta M^{\langle 1 \rangle} + E^{\langle 1 \rangle}) = \textsf{ct}^{\langle 1 \rangle} = ( A_0^{\langle 1 \rangle}, A_1^{\langle 1 \rangle}, \ldots, A_{k-1}^{\langle 1 \rangle}, B^{\langle 1 \rangle}) \in \mathcal{R}_{\langle n,q \rangle}^{k + 1}$

$\textsf{GLWE}_{S, \sigma}(\Delta M^{\langle 2 \rangle} + E^{\langle 2 \rangle}) = \textsf{ct}^{\langle 2 \rangle} = ( A_0^{\langle 2 \rangle}, A_1^{\langle 2 \rangle}, \ldots, A_{k-1}^{\langle 2 \rangle}, B^{\langle 2 \rangle}) \in \mathcal{R}_{\langle n,q \rangle}^{k + 1}$

$ $

\noindent 定义如下密文加法运算：

$\textsf{ct}^{\langle 1 \rangle} + \textsf{ct}^{\langle 2 \rangle} = ( A_0^{\langle 1 \rangle} + A_0^{\langle 2 \rangle}, \text{ } A_1^{\langle 1 \rangle} + A_1^{\langle 2 \rangle}, \ldots, A_{k-1}^{\langle 1 \rangle} + A_{k-1}^{\langle 2 \rangle}, \text{ } B^{\langle 1 \rangle} + B^{\langle 2 \rangle} )$

$ $

\noindent 则有以下结论：

\begin{tcolorbox}[title={\textbf{\tboxlabel{\ref*{sec:glwe-add-cipher}} GLWE 同态加法}}]
$\textsf{GLWE}_{S, \sigma}(\Delta M^{\langle 1 \rangle} + E^{\langle 1 \rangle} ) + \textsf{GLWE}_{S, \sigma}(\Delta M^{\langle 2 \rangle} + E^{\langle 2 \rangle}) $

$ = ( \{A_i^{\langle 1 \rangle}\}_{i=0}^{k-1}, \text{ } B^{\langle 1 \rangle}) + (\{A_i^{\langle 2 \rangle}\}_{i=0}^{k-1}, \text{ } B^{\langle 2 \rangle}) $

$ = ( \{A_i^{\langle 1 \rangle} + A_i^{\langle 2 \rangle}\}_{i=0}^{k-1}, \text{ } B^{\langle 1 \rangle} + B^{\langle 2 \rangle} ) $

$= \textsf{GLWE}_{S, \sigma}(\Delta(M^{\langle 1 \rangle} + M^{\langle 2 \rangle}) + E^{\langle 3 \rangle})$ \textcolor{red}{ $\rhd$ 其中 $E^{\langle 3 \rangle} = E^{\langle 1  \rangle} + E^{\langle 2 \rangle}$}
\end{tcolorbox}


相位恒等式在代数上无条件成立。若 $\|E^{\langle1\rangle}+E^{\langle2\rangle}\|_\infty<\Delta/2$，解码才会返回 $M^{\langle1\rangle}+M^{\langle2\rangle}\bmod t$；例如，两个输入的噪声界之和小于 $\Delta/2$ 就是充分条件。

$ $

%\noindent \textbf{\underline{Proof}}
\begin{myproof}
\begin{enumerate}
\item 定义以下记号： \\
$A_0^{\langle 3 \rangle} = A_0^{\langle 1 \rangle} + A_0^{\langle 2 \rangle}$ \\
$A_1^{\langle 3 \rangle} = A_1^{\langle 1 \rangle} + A_1^{\langle 2 \rangle}$ \\
$\vdots$ \\
$A_{k-1}^{\langle 3 \rangle} = A_{k-1}^{\langle 1 \rangle} + A_{k-1}^{\langle 2 \rangle}$ \\
$E^{\langle 3 \rangle} = E^{\langle 1 \rangle} + E^{\langle 2 \rangle}$ \\
$B^{\langle 3 \rangle} = B^{\langle 1 \rangle} + B^{\langle 2 \rangle}$
\item 可推导得： \\
$B^{\langle 3 \rangle} = B^{\langle 1 \rangle} + B^{\langle 2 \rangle}$ \\
$ = \sum\limits_{i=0}^{k-1}{(A_i^{\langle 1 \rangle} \cdot S_i)} + \Delta \cdot M^{\langle 1 \rangle} + E^{\langle 1 \rangle} + \sum\limits_{i=0}^{k-1}{(A_i^{\langle 2 \rangle} \cdot S_i)} + \Delta \cdot M^{\langle 2 \rangle} + E^{\langle 2 \rangle}$ \\ 
$= \sum\limits_{i=0}^{k-1}{((A_i^{\langle 1 \rangle} + A_i^{\langle 2 \rangle}) \cdot S_i)} + \Delta \cdot (M^{\langle 1 \rangle} + M^{\langle 2 \rangle}) + (E^{\langle 1 \rangle} + E^{\langle 2 \rangle})$ \textcolor{red}{ $\rhd$ 由交换律与分配律} \\
$= \sum\limits_{i=0}^{k-1}{(A_i^{\langle 3 \rangle} \cdot S_i)} + \Delta \cdot (M^{\langle 1 \rangle} + M^{\langle 2 \rangle}) + E^{\langle 3 \rangle}$ \\
\item 由于 $B^{\langle 3 \rangle} = \sum\limits_{i=0}^{k-1}{(A_i^{\langle 3 \rangle} \cdot S_i)} + \Delta \cdot (M^{\langle 1 \rangle} + M^{\langle 2 \rangle}) + E^{\langle 3 \rangle}$，

因此 $(A_0^{\langle 3 \rangle}, A_1^{\langle 3 \rangle}, A_2^{\langle 3 \rangle}, \ldots, A_{k-1}^{\langle 3 \rangle}, B^{\langle 3 \rangle})$ 构成密文 $\textsf{GLWE}_{S, \sigma}(\Delta \cdot (M^{\langle 1 \rangle} + M^{\langle 2 \rangle}) + E^{\langle 3 \rangle})$。 \\
\item 因此， \\
$\textsf{GLWE}_{S, \sigma}(\Delta M^{\langle 1 \rangle} + E^{\langle 1 \rangle}) + \textsf{GLWE}_{S, \sigma}(\Delta M^{\langle 2 \rangle} + E^{\langle 2 \rangle})$ \\
$ = ( A_0^{\langle 1 \rangle} + A_0^{\langle 2 \rangle}, \text{ } A_1^{\langle 1 \rangle} + A_1^{\langle 2 \rangle},\ldots, A_{k-1}^{\langle 1 \rangle} + A_{k-1}^{\langle 2 \rangle}, \text{ } B^{\langle 1 \rangle} + B^{\langle 2 \rangle} )$ \\
$ = ( A_0^{\langle 3 \rangle}, A_1^{\langle 3 \rangle}, A_2^{\langle 3 \rangle}, \ldots,  A_{k-1}^{\langle 3 \rangle}, B^{\langle 3 \rangle})$ \\
$ = ( \{A_i^{\langle 3 \rangle}\}_{i=0}^{k-1}, B^{\langle 3 \rangle})$ \\
$= \textsf{GLWE}_{S, \sigma}(\Delta (M^{\langle 1 \rangle} + M^{\langle 2 \rangle}) + E^{\langle 3 \rangle})$

%\begin{flushright}
%\qedsymbol{} 
%\end{flushright}

\end{enumerate}
\end{myproof}

\subsection{讨论}
\label{subsubsec:glwe-add-cipher-discuss} 

\para{噪声消除：}使用秘密密钥作减法会精确消去掩码乘积；只有所得误差各系数的绝对值仍小于 $\Delta/2$ 时，舍入才会去除误差。若合并后的误差超出此界，两个有效输入也可能相加得到无法正确解码的输出。

$ $

\para{噪声增长：}新误差为 $E^{\langle3\rangle}=E^{\langle1\rangle}+E^{\langle2\rangle}$。独立零均值误差的方差相加；若各自方差为 $\sigma^2$，则 $h$ 个独立样本之和的方差为 $h\sigma^2$。仅有零均值并不充分。一般有 $\operatorname{Var}(e_1+e_2)=\operatorname{Var}(e_1)+\operatorname{Var}(e_2)+2\operatorname{Cov}(e_1,e_2)$。密文与自身相加得到 $2e$，其方差为 $4\sigma^2$，并非 $2\sigma^2$。确定性的系数界可由三角不等式给出，无需独立性假设。

$ $

\para{$t\nmid q$ 时的明文溢出：}掩码乘积和密文分量始终按模 $q$ 解释。当 $t\mid q$ 时，明文模 $t$ 的回绕与乘以 $\Delta=q/t$ 相容。当 $t\nmid q$ 且 $\Delta=\lfloor q/t\rfloor$ 时，将未约减明文的每个系数写为 $m=m'+vt$，其中 $m'$ 为中心化代表元，则有效误差是 $e-(q\bmod t)v$。在 \autoref{subsubsec:lwe-noise-bound} 的假设下，要求 $|e-(q\bmod t)v|<\Delta/2$。这一条件依赖溢出次数 $v$，并非仅由 $q\gg t$ 保证。最后模 $t$ 的约减也是消息解密的一部分。

 
